% Architecture and design --- section structure from prd.md 3.1, ownership from Table 3.
%
% This chapter OWNS: the dual-path decomposition, the reflex-path membership rule, the
% validator and state-machine layers, the latency budget and the core allocation.
% It does NOT own: the schema and grammar (Master Ch 3, sec:schema -- one fresh sentence and a
% one-page summary here), any measured latency (Ch 5, Exp-2), the model comparison (Master Ch 4;
% only the selected configuration is named here).
%
% Labels: tab:latency-budget is TAKEN in this document by thesis/generated/exp2_latency_budget.tex
% (the measured table Ch 5 will \input). The design budget here is tab:stage-budget.

\chapter{Architecture and design}
\label{chap:architecture}

Section~\ref{sec:safety-problem} reduced the safety problem to three requirements: a wrong
command must be hard to act on, a stopping command must be fast to act on, and the execution of
any command must be bounded at the point where it takes effect. This chapter turns the first two
into a design; the third is a property of the swarm controller and is realised in
Chapter~\ref{chap:implementation}. Section~\ref{sec:dual-path} divides the runtime into two paths
that read one audio stream. Section~\ref{sec:membership-rule} derives the rule deciding which
utterances may take the fast path, and Section~\ref{sec:validation-layers} specifies the three
checks the slow path's output must pass. Sections~\ref{sec:latency-budget}
and~\ref{sec:resource-allocation} then give each stage a latency target and each path its own
processor cores, so that neither path can spend the time of the other. Every figure in this
chapter is a design target, fixed before any measurement; Chapter~\ref{chap:validation} measures
the system against them.

\section{Dual-path decomposition}
\label{sec:dual-path}
The runtime consists of two paths that read the same audio stream and publish to the same command
bus (Figure~\ref{fig:architecture}), the reflex path and the parse path of
Section~\ref{sec:operational-context}. The reflex path is a two-class keyword spotter~\cite{oww}
that listens continuously and, on detecting either stopping phrase, publishes a fixed command
without transcribing anything. The parse path is the general path: a \gls{vad} model~\cite{silero}
segments the stream into utterances, \texttt{whisper.cpp tiny.en}~\cite{whisper,whispercpp} performs
\gls{stt} on each, a fine-tuned \gls{slm} parses the transcript into a structured command under a
decoding grammar~\cite{llamacpp}, and a semantic validator checks the result before it is published.
The two paths share nothing between the audio buffer and the bus. They are not a fast and a slow
setting of one pipeline. The parse path cannot begin transcribing until the endpointer has
confirmed that the operator has stopped speaking, and that wait alone exceeds the 150~ms reflex
budget (Section~\ref{sec:latency-budget}); a stopping command on the parse path therefore misses
the reflex budget however the rest of the path is tuned (Section~\ref{sec:engineering-requirements}).
The decomposition follows from that constraint rather than from a preference for redundancy.

The language model on the parse path is Qwen2.5-0.5B~\cite{qwen25}, fine-tuned and quantised to
Q4\_K\_M. The \emph{M\'emoire de Master} selects this configuration by comparing models and
quantisation levels on the same device; the present document treats the choice as fixed and
designs the runtime around it.

\begin{figure}[htbp]
\centering
\resizebox{\textwidth}{!}{%
\begin{tikzpicture}[
    font=\footnotesize,
    stage/.style={draw, rounded corners=2pt, align=center, minimum height=10mm,
                  text width=18mm, fill=white},
    flow/.style={->, thick},
    tag/.style={font=\footnotesize\bfseries}]
  % Raspberry Pi 5
  \node[stage] (mic) at (0,0)      {Microphone};
  \node[stage] (buf) at (2.5,0)    {Ring buffer\\16~kHz};
  \node[stage] (kws) at (5.3,1.4)  {Keyword\\spotter};
  \node[stage] (vad) at (5.3,-1.4) {Endpointer\\(\acrshort{vad})};
  \node[stage] (stt) at (7.8,-1.4) {\acrshort{stt}\\\texttt{tiny.en}};
  \node[stage] (slm) at (10.3,-1.4) {\acrshort{slm}\\under grammar};
  \node[stage] (val) at (12.8,-1.4) {Validator\\(layer 2)};
  \node[stage, text width=21mm] (bus) at (15.4,0) {Command bus\\UDP, \acrshort{json}};
  % Workstation
  \node[stage, text width=21mm] (fsm) at (15.4,-4.0) {State machine\\(layer 3),\\ordering rule};
  \node[stage] (ctl) at (12.3,-4.0) {Swarm\\controller\\50~Hz};
  \node[stage] (sim) at (9.3,-4.0)  {Simulated\\swarm\\$N = 5$};

  \draw[flow] (mic) -- (buf);
  \draw[flow] (buf.east) -- ++(0.45,0) |- (kws.west);
  \draw[flow] (buf.east) -- ++(0.45,0) |- (vad.west);
  \draw[flow] (vad) -- (stt);
  \draw[flow] (stt) -- (slm);
  \draw[flow] (slm) -- (val);
  \draw[flow] (kws.east) -| node[pos=0.25, above] {fixed command} (bus.north);
  \draw[flow] (val.north) |- (bus.west);
  \draw[->, thick, dashed] (kws.south) -- ++(0,-0.95) -| node[pos=0.25, above] {cancel decode}
      (slm.north);
  \draw[->, thick, dashed] (bus.south) -- node[right] {Wi-Fi} (fsm.north);
  \draw[flow] (fsm) -- (ctl);
  \draw[flow] (ctl) -- (sim);

  \node[tag, anchor=south west] (ta) at (kws.north west) {Reflex path};
  \node[tag, anchor=north west] (tb) at (vad.south west) {Parse path};

  \begin{scope}[on background layer]
    \node[draw, fill=gray!8, rounded corners=4pt, inner sep=9pt,
          fit=(mic) (ta) (tb) (val) (bus),
          label={[font=\footnotesize\itshape, anchor=north west]north west:Raspberry Pi 5, offline}] {};
    \node[draw, fill=gray!8, rounded corners=4pt, inner sep=9pt,
          fit=(sim) (fsm),
          label={[font=\footnotesize\itshape, anchor=north west]north west:Workstation}] {};
  \end{scope}
\end{tikzpicture}}
\caption[Dual-path runtime architecture]{Dual-path runtime architecture. One audio stream feeds two
paths. The reflex path, the keyword spotter, publishes one of two fixed commands and bypasses
transcription and parsing entirely; the parse path endpoints the utterance, transcribes it, parses
the transcript into a command under the decoding grammar and validates the result. Both paths
publish to one command bus. A reflex-path trigger cancels any decode in progress (dashed); the
ordering rule at the bus's consumer discards any parse-path result older than the latest reflex-path
command, whether or not the cancellation succeeded. The state machine, the swarm controller and the
simulated vehicles run on a workstation joined to the Pi by Wi-Fi.}
\label{fig:architecture}
\end{figure}

\paragraph{Placement.} Everything up to and including the command bus runs on the
Raspberry~Pi~5. The flight state machine, the swarm controller and the simulated aircraft run on a
workstation. The split follows the claim under test: this document asks whether perception and
language understanding fit on the edge device, not whether a physics engine does. A simulator
sharing the Pi's four cores would contend for the cores that Section~\ref{sec:resource-allocation}
allocates to the pipeline, and every latency figure would then depend on the simulator's load as
well as on the pipeline. With the simulation off the device, latency measured on the Pi is a
property of the pipeline alone. The state machine is placed beside the controller rather than
beside the parser because two of its transitions, reaching the take-off altitude and making
ground contact, are observations of the vehicles, and only the controller's side of the bus
makes them.

\paragraph{Command bus.} The two machines are joined by a minimal publish--subscribe bus
carrying single-line \gls{json} messages over UDP. Each message carries a command, a tag naming
the path that produced it, and a sequence number drawn from one monotonic counter shared by both
paths. The bus implements nothing else. The rule deciding whether a message may take effect
belongs to its consumer, beside the state machine, and messages from both paths reach the state
machine by that same route; a reflex-path message differs from a parse-path message only in that
it is a constant, never the output of a decoder. The bus sits behind an interface so that a
distributed middleware can replace it in a deployment beyond simulation. Adopting one here would
add failure modes of its own, such as peer discovery and network partition, to a system whose
experiments are not designed to measure them.

\paragraph{Prompt and prefix.} The command schema is not described in the prompt. Fine-tuning
teaches the model the command format, so the prompt carries only a fixed system instruction and
the transcript, and the fixed part is identical for every request. The design keeps that fixed
part resident in the model's key--value cache between requests, so that prefill for each command
processes only the transcript's tokens. The design estimate is roughly 15 tokens per command,
against roughly 290 if the schema had to be restated in every prompt, and it is the reason a
512-token context window suffices. The deployed parser departs from this assumption, for a reason
given in Chapter~\ref{chap:implementation}; the latency budget of Section~\ref{sec:latency-budget}
is stated as designed.

\paragraph{Preemption and ordering.} When the reflex path fires while the parse path is decoding,
two actions follow, and only one of them is needed for correctness. The needed one is an ordering
rule: the consumer discards any parse-path message whose sequence number is lower than that of the
latest reflex-path message already applied. A parse-path command that was being decoded when the
operator said \emph{swarm hold} therefore cannot take effect after the hold, whether or not its
decode was stopped. The rule depends on when a sequence number is assigned. A parse-path utterance
reserves its number when it is accepted for processing, before transcription begins, and not when
its decode completes. A number assigned at completion would make a slow decode that finishes after a
later trigger compare as newer than the trigger, and the stale command would override the stop the
rule exists to protect. The second action is an optimisation: the decode in progress is cancelled,
so that the three inference cores are released at once rather than spent on a result that will be
discarded. Correctness does not depend on the cancellation succeeding. What the cancellation governs
is recovery, the time before the parse path can accept the next utterance, which the
preemption-recovery criterion bounds at 300~ms at p95 (Table~\ref{tab:nonfunctional-requirements}).

\section{Reflex-path membership rule}
\label{sec:membership-rule}
The reflex path bypasses the stages that interpret speech, not the stages that check commands. Its
commands enter the same consumer as the parse path's and pass the semantic validator and the state
machine like any other message; what they bypass is transcription and parsing, the two stages that
establish what the operator said. The consequence is that no layer can detect a false accept. A
spurious trigger publishes \texttt{\{"intent":"hover"\}}, which is well-formed, inside the physical
envelope and legal while the swarm is taking off or flying: it is indistinguishable, at every check
the architecture has, from a hold the operator asked for. The layers of
Section~\ref{sec:validation-layers} test the form, range and legality of a command, and a false
accept on the reflex path fails none of those tests, because the error lies in whether the command
was spoken at all. Because the spotter listens continuously, every class it carries is a standing
source of such errors. The only protection is therefore to restrict which commands the path may
carry, and the membership rule does so with two tests, both of which a command must pass.

\paragraph{Latency-decisive.} The first test asks whether the time the reflex path saves changes the
physical outcome. The difference between the parse budget and the reflex budget is $2{,}500 - 150 =
2{,}350$~ms (Table~\ref{tab:nonfunctional-requirements}). At the validator's speed ceiling of
2.0~m/s, that interval is 4.7~m of travel, comparable to the largest formation radius the envelope
admits (10~m) and larger than every admissible inter-drone spacing (1--5~m). A stop that arrives
through the parse path can therefore arrive after a vehicle has already covered the distance the
stop was meant to prevent. A command passes this test only if that delay matters to what the
vehicles do.

\paragraph{Fail-safe under false accept.} The second test asks what happens when the command is
executed without having been spoken. A command passes only if its unprompted execution moves the
swarm towards a safer state: the reflex path may make the swarm do less, never more. This is the
test that answers the objection of the previous paragraph, since a false accept of a command that
passes it costs an unnecessary stop and nothing else.

Neither test is sufficient alone. The first alone would admit any urgent movement, and a false
accept of an urgent movement is uncommanded motion with no stage behind it able to recognise the
error. The second alone would admit commands that are harmless but gain nothing from the fast
path, and each of them would enlarge the false-accept surface for no benefit, because the
keyword false-accept budget is aggregated over every reflex-path class. Table~\ref{tab:membership}
applies both tests to the ten intents of the schema. Exactly two pass.

\begin{table}[htbp]
\centering
\footnotesize
\caption[Reflex-path membership rule applied to the ten intents]{The reflex-path membership rule
applied to every intent of the command schema. A command is admitted to the reflex path only if it
passes both tests. The Path column names the path that carries the command. A dash marks a test
that was not evaluated because the command had already failed the other; ``Parse only'' marks a
command whose unprompted execution would add motion.}
\label{tab:membership}
\begin{tabularx}{\textwidth}{@{}l X X l@{}}
\toprule
Intent & Latency-decisive & Fail-safe under false accept & Path \\
\midrule
\texttt{hover}     & Yes & Yes: stops lateral motion & Reflex \\
\texttt{abort}     & Yes & Yes: descends in place & Reflex \\
\texttt{land}      & No: the descent itself is slow & Yes & Parse \\
\texttt{set\_param} & No & Yes, for a value that reduces speed or spacing & Parse \\
\texttt{unknown}   & No: nothing is dispatched & Yes & Parse \\
\texttt{takeoff}   & -- & No: spins up the rotors & Parse only \\
\texttt{move}      & -- & No: uncommanded translation & Parse only \\
\texttt{altitude}  & -- & No: uncommanded climb or descent & Parse only \\
\texttt{rotate}    & -- & No: uncommanded yaw & Parse only \\
\texttt{formation} & -- & No: every vehicle changes slot & Parse only \\
\bottomrule
\end{tabularx}
\end{table}

The two admitted classes map onto intents the schema already contains, so the reflex path adds no
grammar rule, no schema field and no controller entry point. \texttt{swarm hold} publishes
\texttt{hover}: lateral motion stops, altitude is held, the rotors keep running, and the mission can
be resumed. The hover is already the state every rejection on the parse path resolves to where the
flight state permits it (Section~\ref{sec:validation-layers}), so the reflex makes the system's
existing safe state reachable at reflex speed rather than introducing a new one. \texttt{swarm
abort} publishes \texttt{abort}, defined as an immediate descent in place with the motors stopped on
ground contact; the mission ends and cannot be resumed by voice. Two other definitions were
rejected. Stopping the motors in flight, anywhere in an altitude envelope of 0.5--15~m, is a
decision to lose the airframe. A return to the launch point requires flight, which is what the
command exists to stop. Descending in place removes energy from the situation without lateral
motion, and it keeps the two classes distinct in effect: a hold stays airborne and an abort comes
down.

Both phrases are two words with a common prefix rather than single words. A keyword classifier
scores a fixed window of audio~\cite{oww}, and a single short word fills little of it, so a
one-word class gives the classifier less evidence and invites more false accepts on continuous
speech. The shared word \emph{swarm} supplies an onset that is uncommon in ordinary speech to
both classes at once. It also brings the two classes closer acoustically, which raises the chance
of one being heard as the other. That confusion is benign by construction: hearing \emph{abort}
for \emph{hold} is the more conservative error, and hearing \emph{hold} for \emph{abort} still
stops the swarm. Because both classes passed the second test, an error between them stays within
the set of fail-safe outcomes.

The most obvious third class is excluded deliberately. A \texttt{resume} phrase would be as
urgent to the operator as any other, but a false accept would restart motion the operator had
chosen to freeze, which reverses a safety action rather than adding one. It fails the second test,
and resumption is therefore available only through the parse path, where the added 2.35~s costs
nothing. Recovery out of an abort is not available by voice on either path, for the same reason
(Section~\ref{sec:validation-layers}).

The rule states the condition under which a third class could be admitted: its command passes
both tests, and the false accepts it adds fit within the aggregate budget measured with it in
place. The tests apply to the command a phrase publishes, not to the phrase, so a change of
wording cannot admit a command the rule excludes, and they are applied again whenever the schema
gains an intent. For the present vocabulary the rule yields exactly two classes; that result is
the contribution of this document (Section~\ref{sec:contributions}), and it is what makes the
membership of the reflex path checkable by a reader who did not choose it.

\section{Three validation layers}
\label{sec:validation-layers}
% - Layer 1 (one page, write-once): one fresh sentence naming the Memoire de Master as owner of
%   the schema and grammar; then only what this document needs from it -- ten intents, bounded
%   digits and identifier list, structure guaranteed, ranges not expressible. No grammar listing.
% - Layer 2, semantic validator: required-slot matrix; envelope clamp-and-log, never reject
%   (radius [1,10] m, spacing [1,5] m, z [0.5,15] m, speed [0.2,2.0] m/s, |pos| <= 50 m by
%   vector scaling so direction survives, yaw wrapped into [-180,180] deg); identifiers filtered
%   to {0..N-1}, deduplicated, sorted, empty -> all drones; any failure -> HOVER + log.
%   Source of record: schema/schema.py ENVELOPE, schema/validate.py, ADR-0001.
% - Layer 3, flight state machine: five states, transitions, Table tab:legality (Table 9).
%   Rejection resolution per ADR-0002: HOVER where hover is legal (TAKING_OFF, FLYING), logged
%   no-op where it is not (LANDED, LANDING, ABORTED) -- the one place the "every rejection is a
%   hold" principle of Ch 1 needs its precise form. abort in LANDED/ABORTED is a logged no-op.
%   Recovery out of ABORTED is manual and non-vocal: keeps the reflex one-directional.
% - What no layer catches: a schema-valid, in-envelope, state-legal wrong command (the yaw sign
%   flip). Named as a boundary of the design, measured in Ch 5 / discussed in Ch 6. No figure.
\TODO{draft}

\section{Latency budget}
\label{sec:latency-budget}
% Budget = DESIGN. No measured figure anywhere in this section; Chapter~\ref{chap:validation}.
% - Table tab:stage-budget: both paths, both anchors -- endpointing wait, STT, SLM prefill,
%   SLM decode, validate+FSM+dispatch, E2E from end of speech, E2E from start of speech (3 s
%   utterance), reflex from keyword offset, reflex from keyword onset (~700 ms phrase).
%   Row wording fresh (Master's tab:latency-budget restates the parse rows).
% - Anchors: T0 = last frame in the buffer when the endpointer declares end of speech; reflex
%   anchored at keyword offset because an onset target of 150 ms is unsatisfiable by
%   construction; onset reported as operator-perceived time.
% - Endpointing line derived, not tuned: 450 ms of silence plus at most one 32 ms window
%   = 482 ms worst case, inside the 500 ms allowance.
% - Stage allowances are per-stage targets, not a decomposition: the E2E target is judged on the
%   measured E2E, not the sum.
% - Outside the budget, stated so it is not mistaken for an omission: capture/buffering latency
%   before T0, and the Wi-Fi hop after dispatch.
% - One fresh sentence: the Memoire de Master uses the parse-path rows as its deployment
%   constraint.
\TODO{draft}

\section{Resource allocation}
\label{sec:resource-allocation}
% - Table tab:core-allocation (Table 7): core 0 = capture, ring buffer, endpointer, keyword
%   spotter (the 80 ms frame loop); cores 1--3 = whisper.cpp -t 3 and llama.cpp -t 3.
% - Why three cores suffice for both: STT and SLM are sequential by construction.
% - Why the frame loop must never wait on the parse path: reflex latency holds under load only if a decode
%   cannot delay the next frame -- the design property the loaded condition of Exp-2 tests.
%   That isolation held is Ch 5's to report; forward reference only.
% - Per-thread pinning; a child process inherits the spawning thread's mask.
% - Memory: worst-case arithmetic (largest configuration, full stack) as a bound, not a
%   measurement -- OPEN QUESTION for the author: include here, or leave memory to the Master's
%   requirements table?
\TODO{draft}
